\chapter{Otros aspectos de la Ingeniería del Software}

\section{Planificación y gestión de proyectos software}

Gestionar un proyecto software consiste en conseguir que se entregue en plazo,
dentro del presupuesto y con la calidad acordada. Es una actividad de ingeniería
tanto como el diseño, y comparte con él la dificultad de decidir con información
incompleta.

Los proyectos software tienen rasgos que los distinguen de otros: el producto es
intangible, así que el avance no se ve; no hay dos proyectos iguales, de modo que
la experiencia previa se traslada mal; y el proceso cambia con la tecnología a un
ritmo que envejece rápido los datos históricos \cite{sommerville2016}.

\subsection*{Estimación}

Estimar es predecir esfuerzo, plazo y coste. Los enfoques habituales son tres, y
suelen combinarse:

\begin{itemize}
    \item \textbf{Juicio de expertos}, con varias personas estimando por separado
          y contrastando después. Es barato y sorprendentemente robusto cuando el
          equipo conoce el dominio.
    \item \textbf{Estimación por analogía}, comparando con proyectos anteriores
          medidos.
    \item \textbf{Modelos algorítmicos}, como COCOMO, que calculan el esfuerzo a
          partir de una medida del tamaño y de factores de ajuste.
\end{itemize}

Todos comparten el mismo punto débil: dependen de una estimación del tamaño que se
hace cuando menos se sabe del sistema. De ahí la práctica de estimar por rangos,
revisar la estimación en cada iteración y desconfiar de las cifras cerradas dadas
antes de la especificación.

Conviene recordar además que el esfuerzo y el plazo no son intercambiables.
Añadir personas a un proyecto retrasado aumenta el coste de coordinación y suele
retrasarlo más, observación que Brooks formuló hace medio siglo y que sigue
describiendo bien lo que ocurre.

\subsection*{Planificación y riesgos}

La planificación descompone el trabajo en tareas con dependencias, asigna
recursos y fija hitos. Los hitos deben corresponder a entregables verificables;
un hito que se declara cumplido porque «está casi terminado» no informa de nada.

La gestión de riesgos es la parte de la planificación que se salta con más
frecuencia y la que más pronto se echa en falta. Consiste en identificar lo que
puede salir mal, estimar su probabilidad y su impacto, decidir qué hacer con cada
riesgo —evitarlo, mitigarlo, transferirlo o aceptarlo— y revisar la lista de forma
periódica. Los riesgos habituales en software son la rotación del personal, el
cambio de requisitos, la subestimación del tamaño y la dependencia de tecnología
o proveedores no probados.

\subsection*{Calidad y métricas}

La gestión de la calidad establece qué se entiende por calidad en el proyecto y
cómo se comprueba: estándares de producto, estándares de proceso, revisiones e
inspecciones. Las métricas —tamaño, complejidad, densidad de defectos, cobertura
de pruebas— sirven para tomar decisiones, no para evaluar personas. Usadas como
objetivo, dejan de medir lo que medían, porque el equipo optimiza el indicador y
no la propiedad que representaba.

\section{Validación y verificación de software}

Verificar es comprobar que el sistema se construye conforme a su especificación;
validar es comprobar que la especificación describe lo que el cliente necesita. La
distinción clásica lo resume bien: la verificación pregunta si se está
construyendo el producto correctamente, y la validación, si se está construyendo
el producto correcto.

\subsection*{Pruebas}

Las pruebas ejecutan el sistema con la intención de encontrar defectos. Un
resultado sin fallos no demuestra que no los haya: las pruebas muestran la
presencia de errores, nunca su ausencia. Se organizan en niveles:

\begin{itemize}
    \item \textbf{Unitarias}, sobre componentes aislados, escritas normalmente por
          quien los implementa y automatizadas.
    \item \textbf{De integración}, sobre las interacciones entre componentes, donde
          aparecen los fallos de interfaz que las unitarias no ven.
    \item \textbf{De sistema}, sobre el producto completo, incluidos los requisitos
          no funcionales: carga, seguridad, recuperación.
    \item \textbf{De aceptación}, ejecutadas con el cliente o en sus términos, y
          que pertenecen a la validación más que a la verificación.
\end{itemize}

El diseño de los casos de prueba puede partir de la especificación —prueba de caja
negra, con técnicas como las clases de equivalencia y el análisis de valores
límite— o de la estructura interna —prueba de caja blanca, guiada por criterios de
cobertura—. Ninguno de los dos basta por sí solo: la caja negra no llega a los
caminos no previstos y la caja blanca no detecta lo que falta implementar
\cite{pfleeger2002}.

La \textbf{prueba de regresión} es la que se vuelve a pasar tras cada cambio para
comprobar que no se ha roto lo que funcionaba. Solo es viable si está
automatizada, y por eso la automatización de pruebas es una decisión de proceso
más que una preferencia técnica. En la programación extrema se lleva al extremo
escribiendo la prueba antes que el código \cite{beck2002}.

\subsection*{Revisiones e inspecciones}

No toda verificación requiere ejecutar. Las revisiones e inspecciones examinan
documentos y código con una lista de comprobación, y encuentran defectos que las
pruebas no alcanzan, entre ellos los de la propia especificación. Su ventaja es
que actúan antes: un defecto detectado en la revisión de requisitos cuesta una
fracción de lo que cuesta corregido en producción.

\section{Mantenimiento de software}

El mantenimiento es todo el trabajo que se hace sobre el sistema después de la
entrega, y es la fase más larga y más cara del ciclo de vida. Se clasifica en
cuatro tipos:

\begin{itemize}
    \item \textbf{Correctivo}, para reparar defectos.
    \item \textbf{Adaptativo}, para acomodar cambios del entorno: sistema
          operativo, normativa, sistemas con los que se integra.
    \item \textbf{Perfectivo}, para añadir funcionalidad o mejorar el rendimiento a
          petición del usuario.
    \item \textbf{Preventivo}, para facilitar el mantenimiento futuro sin cambiar
          el comportamiento observable.
\end{itemize}

Contra la intuición, el correctivo es la parte menor. La mayoría del esfuerzo se
va en perfectivo y adaptativo, es decir, en hacer evolucionar un sistema que
funciona \cite{sommerville2016}.

\subsection*{Evolución y deterioro}

Las leyes de Lehman describen lo que le ocurre a un sistema en uso: si no cambia,
deja de resultar útil, porque el entorno sí cambia; y a medida que cambia, su
complejidad crece salvo que se dedique trabajo explícito a contenerla. De ahí que
el mantenimiento preventivo —refactorizar, actualizar dependencias, recuperar la
documentación— no sea un lujo, sino la condición para que los otros tres tipos
sigan siendo asequibles.

La \textbf{deuda técnica} nombra ese fenómeno desde el lado económico: cada
decisión tomada por rapidez en lugar de por calidad genera un interés que se paga
en cada cambio posterior. Es legítimo contraerla de forma consciente y con un plan
para devolverla; el problema aparece cuando se contrae sin registrarla, porque
entonces nadie decide cuándo se paga.

\subsection*{Reingeniería}

Cuando el coste de mantener supera al de rehacer, cabe la reingeniería: reconstruir
el sistema conservando su funcionalidad. Es menos arriesgada que un desarrollo
desde cero, porque el sistema existente documenta el comportamiento esperado,
incluido el que nadie llegó a escribir en ningún requisito. La alternativa de
reescribir de cero pierde ese conocimiento y suele reintroducir defectos que ya
estaban resueltos.
